% methodology.tex -- Data Availability Committee Paper (RU-V6)

\section{Methodology}

\subsection{Protocol Overview}

Shamir-DAC operates in four phases: \emph{share generation}, \emph{share
distribution}, \emph{attestation}, and \emph{recovery} (on-demand). We
describe each phase and present the attestation pipeline algorithm.

\subsection{Share Generation}

Given a batch data blob $D$ of $|D|$ bytes, the operator:

\begin{enumerate}
\item Packs $D$ into field elements $\{d_1, \ldots, d_m\}$ where
  $m = \lceil |D| / 31 \rceil$ and each $d_i \in \mathrm{GF}(p)$ with
  $p$ the BN128 scalar field prime ($p = 2^{254} - \delta$, a 254-bit
  prime). Each field element encodes 31 bytes of data (not 32, to ensure
  all values are $< p$).
\item For each element $d_i$, constructs a random polynomial
  $f_i(x) = d_i + a_{i,1} x + \cdots + a_{i,k-1} x^{k-1}$ over
  $\mathrm{GF}(p)$ with uniformly random coefficients
  $a_{i,j} \stackrel{\$}{\leftarrow} \mathrm{GF}(p)$.
\item Evaluates $f_i(j)$ for each committee member $j \in \{1, \ldots, n\}$,
  producing share $s_{i,j} = f_i(j)$.
\item Assembles member $j$'s share bundle:
  $\mathbf{s}_j = (s_{1,j}, s_{2,j}, \ldots, s_{m,j})$.
\end{enumerate}

The computational cost is $O(k \cdot n)$ field multiplications per element,
or $O(k \cdot n \cdot m)$ total. For $(2,3)$-Shamir with a 500~KB batch
($m = 16{,}130$ elements), this is $\approx 96{,}780$ field multiplications.

\subsection{Share Distribution and Attestation}

Algorithm~\ref{alg:attestation} presents the attestation pipeline.

\begin{algorithm}[h]
\caption{Shamir-DAC Attestation Pipeline}
\label{alg:attestation}
\begin{algorithmic}[1]
\Require Batch data $D$, committee $\{C_1, \ldots, C_n\}$, thresholds $(k, m)$
\Ensure Attestation certificate or fallback to on-chain DA
\State $h \gets \text{SHA-256}(D)$ \Comment{Data commitment}
\State $\{d_1, \ldots, d_m\} \gets \text{Pack}(D, 31)$ \Comment{Field elements}
\For{$i = 1$ to $m$}
  \State $f_i(x) \gets d_i + \sum_{j=1}^{k-1} a_{i,j} x^j$ with $a_{i,j} \stackrel{\$}{\leftarrow} \mathrm{GF}(p)$
\EndFor
\For{$j = 1$ to $n$}
  \State $\mathbf{s}_j \gets (f_1(j), f_2(j), \ldots, f_m(j))$
  \State Send $(\mathbf{s}_j, h)$ to $C_j$
\EndFor
\State $\text{sigs} \gets \emptyset$
\For{$j = 1$ to $n$}
  \State $C_j$ stores $\mathbf{s}_j$, verifies $h$, returns $\sigma_j \gets \text{ECDSA.Sign}(sk_j, h)$
  \State $\text{sigs} \gets \text{sigs} \cup \{(j, \sigma_j)\}$
\EndFor
\If{$|\text{sigs}| \geq m$}
  \State \Return \text{Certificate}$(h, \text{sigs}, \text{timestamp})$
\Else
  \State Post $D$ on-chain \Comment{AnyTrust fallback}
\EndIf
\end{algorithmic}
\end{algorithm}

The certificate is submitted to the L1 smart contract
(\texttt{DACAttestation.sol}), which verifies $m$ ECDSA signatures via
\texttt{ecrecover}. On Basis Network (zero-fee Subnet-EVM), the verification
cost of $3 \times 3{,}000 = 9{,}000$ gas is computationally negligible.

\subsection{Data Recovery}

When reconstruction is needed (e.g., an enterprise requests a withdrawal proof
or an auditor requires data verification), any $k$ committee members provide
their share bundles. For each element index $i$, the secret $d_i$ is recovered
via Lagrange interpolation:

\begin{equation}
d_i = f_i(0) = \sum_{j \in S} s_{i,j} \cdot \prod_{\ell \in S, \ell \neq j}
\frac{\ell}{\ell - j} \mod p
\label{eq:lagrange}
\end{equation}

\noindent where $S$ is any subset of $k$ member indices. The recovered
elements are concatenated and unpacked to obtain the original byte sequence
$D$. Reconstruction cost is $O(k^2)$ field multiplications per element.

\subsection{Information-Theoretic Privacy Argument}

The privacy guarantee follows directly from the information-theoretic security
of Shamir's scheme~\cite{shamir1979share}. For a $(k,n)$-threshold scheme
over $\mathrm{GF}(p)$:

\begin{itemize}
\item Any $k-1$ shares are jointly uniformly distributed over
  $\mathrm{GF}(p)^{k-1}$, independent of the secret.
\item The mutual information $I(S; \mathbf{s}_{T}) = 0$ for any secret $S$
  and any set $T$ of shares with $|T| < k$.
\item This holds against computationally \emph{unbounded} adversaries---no
  cryptographic assumption is required.
\end{itemize}

For enterprise batch data packed into $m$ field elements, each element is
independently shared. The information leakage from $k-1$ share bundles
about the entire batch is:

\begin{equation}
I(D; \mathbf{s}_{T_1}, \ldots, \mathbf{s}_{T_{k-1}}) =
\sum_{i=1}^{m} I(d_i; s_{i,T_1}, \ldots, s_{i,T_{k-1}}) = 0
\label{eq:privacy}
\end{equation}

\noindent This is strictly stronger than computational privacy (encryption),
which requires assumptions about adversary capabilities. Contrast with
EigenDA, where each chunk contains partial plaintext data, and Semi-AVID-PR,
where privacy relies on the computational hardness of breaking the encryption
scheme~\cite{nazirkhanova2022semiavidpr}.

\subsection{Comparison with Alternative Approaches}

Table~\ref{tab:approaches} compares the design alternatives for private
data availability in validium systems.

\begin{table}[h]
\centering
\caption{Design Alternatives for Private Data Availability}
\label{tab:approaches}
\begin{tabular}{lp{1.4cm}p{1.6cm}p{1.6cm}}
\toprule
\textbf{Approach} & \textbf{Privacy} & \textbf{Storage Overhead} & \textbf{Recovery Cost} \\
\midrule
Full replication & None & $n \times$ & $O(1)$ read \\
Reed-Solomon     & None (partial plaintext) & $n/k \times$ & $O(k \log^2 k)$ \\
Encrypt-then-RS  & Computational & $n/k \times$ & $O(k \log^2 k)$ + decrypt \\
Semi-AVID-PR     & Computational & $\sim$3.2$\times$ & $O(n/3)$ nodes \\
\textbf{Shamir-DAC} & \textbf{Info.-theoretic} & $n \times$ & $O(k^2 \cdot m)$ \\
\bottomrule
\end{tabular}
\end{table}

Shamir-DAC trades higher per-node storage ($1\times$ original per member,
$n\times$ total) for the strongest possible privacy guarantee. For enterprise
batch sizes ($\leq$1~MB) and small committees ($n \leq 5$), the storage
overhead is negligible (e.g., 3.87~MB total for a 1~MB batch across 3 nodes).
The trade-off favors Shamir when: (a) batch sizes are small, (b) committee
sizes are small, and (c) information-theoretic privacy is required.
